\chapter{Eye Tracker}
\label{sec:eye tracker}

In diesem Kapitel wird ein grober technischer Einblick in die Funktionen eines Eye Trackers bereitgestellt. Daraufhin wird eine Marktanalyse durchgeführt, um verfügbare Eye Tracker auf deren Anwendbarkeit, Integration in CARLA und Kosten zu untersucht.

\section{Technische Grundlagen}

Die Hardwarevoraussetzungen für einen Eye-Tracker belaufen sich auf eine Kamera, eine Lichtquelle und einen Computer. Als Kamera genügt in einigen Anwendungen eine einfache Webcam. Die Genauigkeit des Eye Trackers kann gesteigert werden, indem eine hochauflösende Kamera genutzt wird. Die Auflösung hängt allerdings auch stark mit dem Abstand der Testperson zur Kamera ab. Sitzt die Testperson recht nahe an der Kamera, dann ist auch eine kleinere Auflösung ausreichend. In unseren Experimenten wurde eine Kamera mit einer Auflösung von 1920x1080 Pixel eingesetzt. Diese Auflösung hat sich in Kombination mit einem Abstand von unter einem Meter zur Kamera als ausreichend erwiesen. Auch Kameras mit höherer Bildrate können bessere Ergebnisse liefern. Allerdings reicht eine Bildwiederholrate von 30 Bildern pro Sekunde für die Erfassung von Augenbewegungen vollkommen aus  \cite{lit_tracking_framerate}. Kameras haben im Vergleich zum menschlichen Auge die Fähigkeit infrarotes Licht zu erfassen. Digitalkameras und Handykameras können je nach eingesetzter Sensortechnologie bereits IR-Strahlung wahrnehmen. Oft sind allerdings Filter eingebaut, die IR-Strahlung größtenteils herausfiltern, da die Kameras die menschliche Sicht abbilden sollen. Es gibt allerdings spezielle IR-Kameras, deren Sensor auf einer anderen Technologie basiert, sodass diese besonders im IR-Bereich sensitiv sind. Einsatzgebiete sind bspw. Wärmebildkameras oder Nachtsichtkameras. Im Kontext von Eye Tracking Systemen können solche Kameras selbst bei schlechten Lichtverhältnissen noch ausreichend Bildinformation erhalten im Vergleich zu einer Kamera, welche nur im menschlich sichtbaren Frequenzbereich operiert. \cite{lit_article_infrared}
Durch das Anbringen von Infrarotlampen und dem Einsatz einer solchen Infrarotkamera können auch bei schlechteren Lichtverhältnissen gute Ergebnisse erzielt werden \cite{lit_article_stateoftheart_eyetracker}.

Auf die genaue Funktionsweise von Eye Tracking Systemen soll in dieser Projektarbeit nicht eingegangen werden, da es nicht Zielsetzung ist ein solches System selbst zu entwickeln. Wichtig ist lediglich, dass die Ergebnisqualität des Eye Tracking System stark von den Lichtverhältnissen und vom Abstand der Testperson zur Kamera abhängig ist. Diese Eigenschaften gelten universal und sind nicht an ein bestimmtes Eye Tracking System gebunden.

\section{Marktanalyse} 

Da die Entwicklung eines eigenen Eye-Tracking Algorithmus den zeitlichen Rahmen der Projektarbeit übersteigt, wird eine Marktanalyse von bereits verfügbaren Eye Trackern durchgeführt.

Ein erster Softwarekandidat ist ein Open Source Tracking Projekt auf GitHub \cite{lit_code_github_facegazetracker}. Die Ausgabe dieses Trackers ist in Abbildung \ref{fig:old_tracker_face_gaze_webcam} zu erkennen. Das System orientiert sich hauptsächlich an den Augen und der Nase der Testperson. Dadurch kann es recht genau bestimmen, in welche Richtung der Kopf einer betrachteten Person geneigt ist. Zudem kann die Software das Blinzeln der Augen erkennen.

Die Software schätzt nicht die Blickposition der Testperson auf dem Bildschirm, sondern gibt nur an, in welcher Lage sich der Kopf befindet. Diese Daten könnten aber als Grundlage für eine Weiterentwicklung dienen. Da allerdings nur die Ausrichtung des Gesichts und nicht der Pupillen in die Ausgabe einfließt, erwarten wir keine besonders genauen Ergebnisse durch diese Herangehensweise.

\begin{figure}[h!]
	\begin{center}
		\includegraphics[scale=0.4]{old_tracker_face_gaze_webcam}
		\caption{Ausgabe des Face-Gaze-Trackers, Quelle: \cite{lit_code_github_facegazetracker}} \label{fig:old_tracker_face_gaze_webcam}
	\end{center}
\end{figure}

Eine weitere Option stellt das open source Projekt "Eyegaze Tracker" dar. Diese Software erfasst nicht die Gesichtsausrichtung, sondern die Pupillen. Dadurch kann sie Aussagen über die grobe Blickrichtung treffen. So wie in Abbildung \ref{fig:old_tracker_eyegaze_webcam} zu erkennen ist, kann bspw. die Aussage getroffen werden, dass die Testperson nach links schaut. Auch hier wird keine Blickposition auf dem Bildschirm geschätzt. Allerdings werden die Pupillen erfasst und die grobe Blickrichtung bestimmt. Die Blickposition kann gegebenenfalls mit den vorhandenen Daten durch eine Weiterentwicklung des vorhandenen open source Codes bestimmt werden. Dies würde allerdings sehr viel Entwicklungsaufwand bedeuten. 

\begin{figure}[h!]
	\begin{center}
		\includegraphics[scale=0.4]{old_tracker_eyegaze_webcam}
		\caption{Ausgabe des Eyegaze-Trackers, Quelle: \cite{lit_code_github_gazetracking}} \label{fig:old_tracker_eyegaze_webcam}
	\end{center}
\end{figure}

Neben den oben genannten Projekten wurden auch noch einige weitere auf ihre Einsetzbarkeit in diesem Projekt untersucht. Diese hatten größtenteils ähnliche Eigenschaften und Einschränkungen. Eye Tracker die einen größeren Funktionsumfang bieten wie bspw. das System "Tobii Eye Tracker" haben den Nachteil, dass sie sehr kostenintensiv sind. Es wurde sich nach viel Recherche und Experimenten für den Beam Eye Tracker von Eyeware Tech entschieden. Dieser hat unsere Anforderungen grundlegend erfüllt und soll als Basis für die weitere Entwicklung dienen.

\section{Beam Eye Tracker}

Der Beam Eye Tracker ist für Gaming-Anwendungen und dem Einsatz von einfacher Hardware wie normalen Webcams entwickelt. Neben der Schätzung der Blickposition auf einen Punkt auf dem Bildschirm verfügt das Programm über eine Kopfpositionserkennung. Durch die Funktion "Gaming Erweiterungen" teilt der Beam Eye Tracker die Trackingdaten mit anderen Programmen. Diese API ermöglicht die Nutzung der Trackingdaten in eigenen Programmen. 

Die Beam API kann über die Programmiersprachen C++, C, C\# und Python angesprochen werden \cite{lit_code_beam_sdk}. Im Rahmen dieser Projektarbeit wird die API mit Pyhton angesprochen. Verschiedene Klassen in der API ermöglichen eine individuelle Einbindung der Eye Tracking Daten in eine selbst entwickelte Anwendung. Relevant für diese Projektarbeit ist die Klasse \glqq{}screen\_gaze\_info\grqq{}. Diese beinhaltet die x- und y- Koordinaten und die Angaben, ob die Augen erkannt werden und wie gut die Schätzung ist.

Durch die einfach einsetzbare API "Beam SDK" und dem niedrigen Preis von 29,99€ des Eye-Trackers fiel die Entscheidung, diesen für das Projekt zu nutzen. In Rücksprache mit den Beam Entwicklern wurde der Beam Eye Tracker für dieses Projekt auf Lebenszeit kostenlos zur Verfügung gestellt.